
\section{Introduction}

The warehouse considered in this thesis is a discrete-event system in which
missions compete for a finite set of mobile resources. A mission is defined
by an origin and a destination, whereas a resource is characterized by its
current position, its compatibility with the mission, and the amount of work
already assigned to it. The scheduling objective is to complete the mission
list within the available working time while avoiding an excessive
concentration of work on a single resource.

This problem has two important characteristics. First, the decisions are
discrete: a mission is assigned to one resource or it is not assigned to that
resource. Second, the duration of an assignment is state-dependent, because
it depends on the current position of the resource and on the path found in
the warehouse map. Consequently, the assignment problem cannot be described
adequately by a purely static sorting rule based only on the mission length.

The approach developed here combines three models with different roles:
\begin{itemize}
	\item the colored Petri net (CPN) represents the logical evolution of the
	mission system and verifies resource availability;
	\item the A* planner verifies geometric reachability and estimates the
	duration of a mission for each candidate resource;
	\item the scheduling MPC ranks the feasible assignments according to a
	prediction of the future resource workload.
\end{itemize}

The resulting controller is referred to as a \emph{discrete scheduling MPC}.
It follows the receding-horizon principle of MPC, but it does not claim to be
a conventional continuous linear MPC or a complete mixed-integer quadratic
program. The distinction is important: the CPN and the path planner enforce
feasibility, while the MPC cost function chooses among the feasible
alternatives.

\section{Relation with Discrete-Event MPC}

In conventional MPC, a model predicts the future evolution of a process over
a prediction horizon. At decision step $k$, a sequence of future inputs is
selected by minimizing a cost function. Only the first input is applied; at
the following decision step, the state is updated and the optimization is
repeated. This is the receding-horizon principle.

The same idea can be used for discrete-event systems. In a warehouse, an
event may correspond to the assignment of a mission, the start of a
transportation activity, or the completion of a mission. Event instants are
not necessarily equally spaced, and the available resources impose logical
constraints. This is consistent with the distinction made in the reference
literature between time-driven systems and event-driven systems.

For max-plus-linear systems, a discrete-event model can be written as
\begin{equation}
	x(k+1) = A \otimes x(k) \oplus B \otimes u(k),
\end{equation}
where $\otimes$ denotes addition and $\oplus$ denotes maximization. Such a
model is useful when the state variables represent event occurrence times
and the system is dominated by synchronization operations.

The present implementation uses a simpler workload model in conventional
algebra. This choice is deliberate. The aim is to obtain a computationally
light assignment mechanism that can be integrated with the existing CPN and
A* code. The workload model predicts accumulated resource times, while the
CPN continues to describe the detailed event logic. A complete max-plus
model of the whole warehouse would require a reformulation of the timed
transitions, resource conflicts, and path-dependent durations.

\section{Scheduling State and Decision Variables}

Let $n_r$ be the number of available resources. At scheduling step $k$, the
workload state is defined as
\begin{equation}
	x(k) =
	\begin{bmatrix}
		T_1(k) & T_2(k) & \cdots & T_{n_r}(k)
	\end{bmatrix}^{\mathsf{T}},
\end{equation}
where $T_r(k)$ is the accumulated working time of resource $r$. In the
MATLAB implementation, this vector corresponds to the field
\texttt{batch.Time}.

The state is complemented by information that is not represented directly in
the workload vector:
\begin{equation}
	z(k) = \left(m(k),p_1(k),\ldots,p_{n_r}(k),\mathcal{M}(k)\right),
\end{equation}
where $m(k)$ is the CPN marking, $p_r(k)$ is the position of resource $r$,
and $\mathcal{M}(k)$ is the set of pending missions. This extended state is
used by the CPN and by the A* planner. The MPC cost function uses primarily
$x(k)$, while feasibility depends on the complete state $z(k)$.

For a mission $i$ and resource $r$, define the binary assignment variable
\begin{equation}
	u_{ir}(k) \in \{0,1\},
\end{equation}
where $u_{ir}(k)=1$ means that mission $i$ is assigned to resource $r$.
For a single decision, the assignment must satisfy
\begin{equation}
	\sum_{r=1}^{n_r} u_{ir}(k) = 1,
\end{equation}
provided that mission $i$ is executed at step $k$. In practice, the CPN
constructs the admissible assignment set before the MPC evaluates the cost,
so the binary constraints are enforced through the color matrices and the
marking rather than through a separate integer optimizer.

\section{CPN-Based Feasibility}

The color matrix associated with a mission identifies which resources are
compatible with it. The marking of the common resource place identifies
which of these resources are currently available. Therefore, a candidate
resource must satisfy both logical and operational conditions.

Let $\mathcal{R}_{\mathrm{CPN}}(i,k)$ denote the resources allowed by the CPN
for mission $i$ at step $k$. This set can contain several resources, because
a mission may be executable by more than one resource. It can also be empty
if no compatible resource has a token available.

The CPN does not calculate the geometric duration of a mission. It only
ensures that the selected resource is represented by a valid color and that
the corresponding token can be consumed. Once a resource is selected, the
CPN updates the marking and advances the current transition. This separation
keeps the logical scheduling model independent from the geometric map model.

In the extended implementation, the CPN provides the set of currently
fireable mission transitions, but it does not impose their execution order.
For every fireable mission, the MPC evaluates the feasible resource choices
and then compares the best mission--resource pair across all fireable
missions. The CPN therefore acts as a feasibility supervisor, while the MPC
also participates in the sequencing decision.

This behavior is enabled by \texttt{selectMissionAndResource = true} in
\texttt{mpcConfig.m}. Setting the flag to \texttt{false} restores selection
of resources in CPN transition order. The comparison is performed among
outgoing transitions of the place currently processed by the simulator;
it is not an optimization across every marked place or across future batches.
Equal ranking keys retain the first candidate in transition order. With
balance priority enabled, the ranking key includes both excess spread and
weighted cost, as detailed below.

The earlier resource-only mechanism evaluated missions in the order in which
their transitions appeared in the CPN. Changing the cost weights could
therefore change the resource assigned to a mission without changing which
mission was considered first. Joint selection expands the decision set: a
later mission can now be selected before an earlier one if its best feasible
assignment has a better ranking key. The CPN still determines the currently
processed place, token availability, and compatible colors. Existing
resource-locking rules also remain part of this admissibility filter.

\section{A* Path Planning and Reachability}

For every candidate resource $r \in \mathcal{R}_{\mathrm{CPN}}(i,k)$, the A*
planner is called with the current resource position, the mission origin,
and the mission destination. The resulting path includes two components:
\begin{enumerate}
	\item the path from the current resource position to the mission origin;
	\item the path from the mission origin to the mission destination.
\end{enumerate}

If the path is feasible, its length is converted into an estimated execution
time. Denoting this time by $d_{ir}(k)$, the duration matrix for a fixed
mission $i$ has one column for each feasible resource candidate $c$:
\begin{equation}
	B_{:,c}(k) = e_{r_c}d_{i r_c}(k),
\end{equation}
where $e_{r_c}$ is the unit vector corresponding to candidate resource $r_c$.
Thus, a column adds time only to its assigned resource. Separate candidate
matrices are evaluated for different missions before their winners are
compared.

A resource for which A* cannot find a path is removed from the candidate set.
The final feasible set is therefore
\begin{equation}
	\mathcal{R}_{\mathrm{feas}}(i,k) =
	\left\{r \in \mathcal{R}_{\mathrm{CPN}}(i,k) \;\middle|\;
	d_{ir}(k) \text{ exists}
	\right\}.
\end{equation}

This filtering is essential. A resource with a low workload must not be
selected if it cannot physically reach the mission. If
$\mathcal{R}_{\mathrm{feas}}(i,k)$ is empty, the mission is deferred and the
simulation records it among the pending missions.

\section{Reduced Workload Model}

The duration matrix is used in the reduced workload model
\begin{equation}
	x(k+1) = A x(k) + B u(k),
\end{equation}
where
\begin{equation}
	A = I_{n_r}.
\end{equation}

The identity matrix represents workload accumulation: the time already
assigned to a resource is retained, and the predicted duration of the new
mission is added to the selected resource. The matrix $B$ represents the
predicted effect of the current assignment.

This model is not intended to reproduce every transition time in the Petri
net. It is an aggregate scheduling model with the following interpretation:
\begin{equation}
	T_r(k+1) = T_r(k) + d_{ir}(k)u_{ir}(k).
\end{equation}

Its advantages are simplicity and low computational cost. Its limitations
are that it does not explicitly model vehicle interactions, congestion,
waiting caused by simultaneous use of map sections, or the full temporal
semantics of the CPN.

\section{Prediction Matrices}

Let $N_p$ denote the prediction horizon. Define the stacked future assignment
vector
\begin{equation}
	U(k) =
	\begin{bmatrix}
		u(k)^{\mathsf{T}} & u(k+1)^{\mathsf{T}} & \cdots &
		u(k+N_p-1)^{\mathsf{T}}
	\end{bmatrix}^{\mathsf{T}}.
\end{equation}
The predicted workload vector is stacked as
\begin{equation}
	Y(k) =
	\begin{bmatrix}
		\hat{x}(k+1|k)^{\mathsf{T}} & \cdots &
		\hat{x}(k+N_p|k)^{\mathsf{T}}
	\end{bmatrix}^{\mathsf{T}}.
\end{equation}

For the linear workload model, the prediction equation is
\begin{equation}
	Y(k) = P x(k) + H U(k),
\end{equation}
where $P$ is the free-response matrix and $H$ is the forced-response
matrix. Their block elements are
\begin{equation}
	P_j = A^j,
\end{equation}
and
\begin{equation}
	H_{j,l} =
	\begin{cases}
		A^{j-l}B, & l \leq j,\\
		0, & l>j.
	\end{cases}
\end{equation}

The matrices are generated in \texttt{mpcPredictionMatrices.m}. The current
configuration uses $N_p=3$. The code evaluates a candidate assignment at the
current step and applies only the first assignment. Thus, the horizon is used
to propagate the workload prediction, while the future assignment entries
are not optimized as an independent multi-step sequence. The method should
therefore be described as a one-step candidate selection embedded in a
receding-horizon prediction framework.

In particular, because $A=I$ and all future inputs after the first are zero,
each predicted block equals $x(k)+Bu(k)$. Increasing $N_p$ alone therefore
does not change the terminal workload or the ranking in this implementation.
The improvements described here come from enlarging the current assignment
set and changing its ranking rule, rather than from optimizing additional
future assignments.

After the chosen assignment is applied, the marking, resource positions, and
workload vector are updated. The prediction is then recomputed using the new
state. This repeated update is the main robustness mechanism of the method.

\section{MPC Cost Criterion}

The reference formulation for conventional MPC uses a quadratic criterion of
the form
\begin{equation}
	J(k) = J_{\mathrm{out}}(k) + \lambda J_{\mathrm{in}}(k),
\end{equation}
with
\begin{equation}
	J(k) =
	\|Y(k)-R(k)\|_Q^2
	+ \lambda\|U(k)\|_R^2.
\end{equation}
For a linear model, this expression can be expanded into a quadratic function
of the future input sequence. In the present problem, however, the future
assignments are binary and the admissible set is defined by the CPN. The
criterion is consequently used to rank a finite set of feasible candidates,
rather than to compute an unconstrained analytical control law.

A direct reference-tracking formulation with an eight-hour reference would
not be appropriate. It could reward a workload close to eight hours even
when a shorter workload is preferable. For this reason, the implementation
uses eight hours as a soft upper bound rather than as a target that must be
tracked from below.

Let
\begin{equation}
	T_{\mathrm{ref}} = 8\,\mathrm{h} = 28800\,\mathrm{s}
\end{equation}
and let $T_r^{\mathrm{pred}}$ be the final predicted workload of resource
$r$. The normalized maximum tardiness is
\begin{equation}
	J_{\mathrm{tardiness}} =
	\frac{\max\left(\max_r T_r^{\mathrm{pred}}-T_{\mathrm{ref}},0\right)}
	{T_{\mathrm{ref}}}.
\end{equation}
This term is zero while all workloads remain below the reference and
increases as soon as the most loaded resource exceeds it.

The workload-balance term is
\begin{equation}
	J_{\mathrm{balance}} =
	\frac{\operatorname{std}(T^{\mathrm{pred}})}{T_{\mathrm{ref}}},
\end{equation}
Here \texttt{std} is the MATLAB sample standard deviation, with normalization
by $n_r-1$ for $n_r>1$; for a single resource the dispersion is zero.
For a candidate resource $c$, the direct load term is
\begin{equation}
	J_{\mathrm{load},c} =
	\frac{T_c^{\mathrm{pred}}}{T_{\mathrm{ref}}}.
\end{equation}
The complete objective used to compare candidates is
\begin{equation}
	J_c =
	w_t J_{\mathrm{tardiness}}
	+ w_b J_{\mathrm{balance}}
	+ w_c J_{\mathrm{load},c}.
\end{equation}

When balance priority is disabled, the selected resource is
\begin{equation}
	c^*(i,k) =
	\underset{c\in\mathcal{R}_{\mathrm{feas}}(i,k)}
	{\operatorname{arg\,min}}\;J_c.
\end{equation}

With joint selection enabled and balance priority disabled, the applied pair is
\begin{equation}
	(i^*,r^*) = \underset{i\in\mathcal{M}_{\mathrm{enabled}}(k),\,
		r\in\mathcal{R}_{\mathrm{feas}}(i,k)}{\operatorname{arg\,min}}\;J_{ir},
\end{equation}
where $\mathcal{M}_{\mathrm{enabled}}(k)$ contains mission transitions
available at the current place. All candidates use the same workload and
resource positions before the winning transition is fired.

\section{Explicit Workload Tolerance and Balance Priority}
\label{sec:balance-priority}

\subsection{Why a larger balance weight is not a tolerance}

The standard-deviation term measures dispersion, but its weighted sum with
the other terms does not enforce a maximum difference between resource
workloads. A smaller value of the weighted cost can compensate for worse
balance. Moreover, below eight hours the tardiness term is identically zero,
so changing its weight alone cannot change the ranking of such candidates.
These properties help explain why tuning weights may have limited effect;
they do not establish the cause of every observed simulation result.

The desired operational quantity is instead the workload range,
\begin{equation}
	\Delta_{ir}=\max_j T_j^{\mathrm{pred}}(i,r)
	-\min_j T_j^{\mathrm{pred}}(i,r),
\end{equation}
computed over all resources in \texttt{batch.Time}, including resources
temporarily unavailable for the current mission. The candidate resource set
and the set used to measure overall balance therefore serve different roles.
The workload is accumulated mission execution time, not a utilization ratio
or a count of assigned missions.

\subsection{Units and excess above the tolerance}

The balance-priority mode, enabled by default with
\texttt{prioritizeBalance = true},
replaces minimization of $J$ alone with lexicographic minimization of
$(v_{ir},J_{ir})$, where
\begin{equation}
	v_{ir}=\max\left(0,\Delta_{ir}-3600\,\delta_h\right).
\end{equation}
Both $\Delta_{ir}$ and $v_{ir}$ are expressed in seconds. The configurable
tolerance \texttt{balanceToleranceHours} is expressed in hours and is initially
$\delta_h=0.02$ hours (72 seconds). The controller first minimizes excess
spread and then the weighted cost, both within each mission and across
missions. If every candidate exceeds the tolerance, it selects the smallest
excess instead of rejecting all assignments. Within tolerance, the weighted
cost decides. Setting \texttt{prioritizeBalance = false} restores the
weighted-cost criterion alone.

The tolerance is an absolute difference between the most and least loaded
resources, not a tolerance of $\pm\delta_h$ around their mean. It is also
independent of the eight-hour reference used to normalize $J$. For example,
workloads of $0.88$ and $0.76$ hours differ by
\begin{equation}
	0.88-0.76=0.12\,\mathrm{h}=432\,\mathrm{s}.
\end{equation}
With a tolerance of $0.02$ hours, their excess is $432-72=360$ seconds.
Workloads of $0.88$ and $0.87$ hours differ by $36$ seconds, giving zero
excess. These numbers illustrate the criterion; they are not a guarantee
that the scheduler can produce either workload vector for a given dataset.

\subsection{Lexicographic comparison}

For two feasible pairs $a$ and $b$, the comparison is
\begin{equation}
	a\prec b
	\quad\Longleftrightarrow\quad
	v_a<v_b\quad\text{or}\quad(v_a=v_b\;\text{and}\;J_a<J_b).
\end{equation}
The applied pair minimizes this ordered key over the feasible set at the
current place:
\begin{equation}
	(i^*,r^*)=
	\underset{(i,r)\in\mathcal{F}(k)}{\operatorname{arg\,min}_{\mathrm{lex}}}
	\left(v_{ir},J_{ir}\right),
\end{equation}
where $\mathcal{F}(k)$ combines CPN and path-planning feasibility.
There are three cases:
\begin{enumerate}
	\item If at least one pair is within tolerance, pairs outside tolerance
	lose to those with $v=0$. The lowest $J$ among the latter is selected.
	\item If all pairs exceed tolerance, the smallest attainable excess is
	selected. Weighted cost decides only among equal excess values.
	\item If both excess and weighted cost are equal, the existing candidate
	order resolves the tie.
\end{enumerate}
This is a strict priority, not an additional weighted penalty. In particular,
a reduction in $J$ cannot compensate for a larger positive excess. The
implementation compares computed excess values directly; it does not round
them into tolerance-sized bins. The tolerance creates a zero-excess region,
within which smaller dispersion can still be rewarded by the existing
standard-deviation term in $J$.

\subsection{Decision scope and practical limits}

This priority can increase travel time or overload relative to minimizing
$J$ alone. It does not guarantee a final spread below tolerance: missions
are indivisible, feasibility may restrict resource choices, and future
missions are not optimized jointly. Final spread is therefore reported
separately from the makespan using actual accumulated workloads.

If a mission is compatible only with a heavily loaded resource, the CPN
filter is respected even when its execution increases the spread. Similarly,
the last available mission may be longer than the remaining imbalance. A
one-step decision cannot split that mission or reassign already completed
work. The priority therefore expresses a desired balance at each decision;
it is not a hard feasibility constraint on every intermediate or final state.

\section{Choice of the Cost Weights}

The weights are not determined by the dimensions of the variables, because
all three objective terms are normalized by $T_{\mathrm{ref}}$. They express
the relative importance assigned to the scheduling objectives within $J$.
The current values in \texttt{mpcConfig.m} are
\begin{equation}
	w_t=0.30,\qquad w_b=0.50,\qquad w_c=0.30.
\end{equation}
Their sum is $1.10$. There is no requirement that the weights sum to one:
dividing all three by $1.10$ would scale $J$ by a positive constant and leave
its ranking unchanged. These weights were retained when balance priority
was introduced. The decisive change is the ordered key $(v,J)$, rather than
another increase in \texttt{balanceWeight}. With priority enabled, weights
affect choices only when their excess values are equal, including all
candidates within tolerance. With priority disabled, they govern the full
comparison through $J$ alone.

This is an engineering initialization, not a claim that the selected triple
is mathematically optimal. A defensible experimental tuning procedure can be
performed by comparing a small set of weight combinations. For example,
consider
\begin{equation}
	\begin{array}{c|ccc}
		\text{configuration} & w_t & w_b & w_c\\
		\hline
		\text{time-oriented} & 0.70 & 0.20 & 0.10\\
		\text{balanced} & 0.60 & 0.25 & 0.15\\
		\text{load-oriented} & 0.50 & 0.30 & 0.20
	\end{array}
\end{equation}

For each configuration, the same Excel mission file, map, initial positions,
and number of resources should be used. The following performance indicators
should be recorded:
\begin{itemize}
	\item final makespan;
	\item maximum resource workload;
	\item standard deviation of resource workloads;
	\item maximum-minus-minimum workload spread and excess over tolerance;
	\item number of completed missions;
	\item number of deferred or unreachable missions;
	\item computational time.
\end{itemize}

A possible selection rule is hierarchical. First discard configurations whose
makespan exceeds $8.5$ hours. Among the remaining configurations, select the
one with the smallest workload standard deviation, using the makespan as a
secondary criterion. This rule is preferable to selecting weights only from
the lowest value of $J$, because the absolute value of $J$ is meaningful
primarily within one fixed weighting configuration.

If the results show that all configurations satisfy the time requirement,
the balanced configuration can be justified by its lower workload dispersion.
If only the time-oriented configuration satisfies the requirement, the
largest weight should be assigned to tardiness. The final choice should be
reported together with the tested alternatives rather than presented as a
universal optimum.

\section{Receding-Horizon Scheduling Procedure}

At each event-driven scheduling step, the following procedure is executed:
\begin{enumerate}
	\item Read the current CPN marking, resource positions, accumulated
	workloads, and pending mission list.
	\item Identify the transitions that can be fired and extract the resource
	candidates allowed by their color matrices.
	\item For every candidate, call A* and calculate the mission duration.
	\item Discard candidates for which the path is infeasible.
	\item Build the duration matrix $B$ and the prediction matrices $P$ and
	$H$.
	\item Evaluate $J_{ir}$ and $v_{ir}$ for each remaining candidate and
	retain the best resource for each fireable mission using the configured
	ranking rule.
	\item Compare the best mission--resource pairs and select the pair with
	the smallest $(v,J)$ key, or the smallest $J$ if priority is disabled.
	\item Convert the selected resource into the color vector required by the
	CPN and fire the transition.
	\item Update the marking, workload, mission history, and resource
	position.
	\item Repeat the procedure with the new state.
\end{enumerate}

The procedure can be summarized by
\begin{equation}
	z(k+1) = \Phi\left(z(k),i^*(k),r^*(k)\right),
\end{equation}
where $\Phi$ denotes the combined update performed by the CPN, the path
planner, and the workload bookkeeping. The MPC does not execute an entire
predicted schedule in advance; it applies only the selected current
assignment and recomputes the decision after the state changes.

\section{MATLAB Implementation}
\label{sec:mpc-matlab-implementation}

The following excerpts show the operative parts of the implementation;
surrounding initialization and path-planning code are omitted where stated.
Standard \texttt{verbatim} environments are used so that the excerpts do not
require a syntax-highlighting package.

\subsection{Configuration and reproducible comparisons}

The relevant settings in \texttt{mpcConfig.m} are:
\begin{small}
	\begin{verbatim}
		options.predictionHorizon = 3;
		options.referenceTime = 8 * 3600;
		options.tardinessWeight = 0.30;
		options.balanceWeight = 0.50;
		options.loadWeight = 0.30;
		options.balanceToleranceHours = 0.02; % 72 seconds
		options.prioritizeBalance = true;
		options.enableDetailedLog = false;
		options.selectMissionAndResource = true;
	\end{verbatim}
\end{small}
The two logical flags control independent aspects of the controller:
\begin{itemize}
	\item \texttt{selectMissionAndResource} controls whether the MPC also
	compares different missions at the current CPN place.
	\item \texttt{prioritizeBalance} controls whether excess spread precedes
	weighted cost in the candidate ranking.
\end{itemize}
To isolate the effect of balance priority, keep joint selection enabled and
change only \texttt{prioritizeBalance}. To reproduce resource selection in
CPN order with weighted cost alone, set both flags to \texttt{false}.
The tolerance remains available for final reporting even when priority is
disabled. It should be nonnegative, and the same tolerance should be used
when comparing results. Reducing it to $0.01$ hours requests a 36-second
range; increasing it to $0.05$ hours permits a three-minute range before
excess becomes positive.

\subsection{Evaluation of an individual assignment}

In \texttt{mpcSelectResource.m}, each candidate activates one entry of the
first input block. The following excerpt evaluates the terminal workload
and its excess:
\begin{small}
	\begin{verbatim}
		assignment(:) = 0;
		assignment(candidate) = 1;
		candidateWorkload = predictedWorkload + H * assignment;
		finalIndex = (options.predictionHorizon-1)*numberOfResources + ...
		(1:numberOfResources);
		finalWorkload = candidateWorkload(finalIndex);
		spreadSeconds = max(finalWorkload) - min(finalWorkload);
		balanceViolations(candidate) = max(0,spreadSeconds - ...
		options.balanceToleranceHours * 3600);
	\end{verbatim}
\end{small}
The function also computes the three normalized terms and their weighted
sum. After evaluating all reachable resources, it selects a candidate using
the shared ranking helper:
\begin{small}
	\begin{verbatim}
		bestCandidate = mpcRankCandidates( ...
		feasibleCosts,balanceViolations,options);
		selectedViolation = balanceViolations(bestCandidate);
	\end{verbatim}
\end{small}
The outputs include the selected resource identifier, its color-matrix row,
all weighted candidate costs, the selected stacked workload prediction, and
the selected excess. The latter is returned separately from the weighted
cost: seconds of excess are not added to a dimensionless objective.

\subsection{Shared ranking rule}

The complete ranking helper \texttt{mpcRankCandidates.m} is reproduced below
with line wrapping adjusted for the page:
\begin{small}
	\begin{verbatim}
		function bestCandidate = mpcRankCandidates( ...
		costs,balanceViolations,options)
		if options.prioritizeBalance
		[~,order] = sortrows( ...
		[balanceViolations(:),costs(:)],[1 2]);
		bestCandidate = order(1);
		else
		[~,bestCandidate] = min(costs);
		end
		end
	\end{verbatim}
\end{small}
The first column is the primary key and the second is the secondary key.
The column conversion \texttt{(:)} ensures compatible shapes regardless of
whether the input arrays are row or column vectors. The helper is called
only after a nonempty feasible candidate set has been obtained.

\subsection{Comparison across missions}

The function \texttt{r\_conflict\_colored.m} retains the winning resource of
each admissible mission. The following excerpt shows the stored quantities:
\begin{small}
	\begin{verbatim}
		candidateTransitions(end+1)=non_zero_idx(i);
		candidateRows(end+1)=resourceRow;
		candidateCosts(end+1)=resourceCosts(resource_index==resourceId);
		candidateResources(end+1)=resourceId;
		candidatePredictions{end+1}=prediction;
		candidateViolations(end+1)=violation;
	\end{verbatim}
\end{small}
The cost must correspond to the resource actually selected under the
$(v,J)$ rule. Using \texttt{min(resourceCosts)} here would be incorrect:
the resource with minimum weighted cost may have lost because of its larger
excess. Mixing its cost with another resource's excess would create a key
that belongs to no evaluated assignment.

The final comparison uses the same helper:
\begin{small}
	\begin{verbatim}
		bestCandidate=mpcRankCandidates( ...
		candidateCosts,candidateViolations,mpcOptions);
		selectedTransition=candidateTransitions(bestCandidate);
		selectedRows=candidateRows(bestCandidate);
		colorIdx={strcat("t",num2str(selectedTransition)),selectedRows};
		right_index=selectedTransition;
		mpcPredictionLog('add',candidatePredictions{bestCandidate});
	\end{verbatim}
\end{small}
Selecting each mission's best pair and then comparing those winners is
equivalent to directly ranking all feasible pairs, because both stages use
the same ordered key. Every discarded resource of a mission is dominated
by that mission's retained resource according to the same comparison.

Only the winning prediction is appended in joint mode. The diagnostic line
reports the mission code, actual resource identifier, weighted cost, and
excess in seconds. A color-matrix row is an index local to an arc; it need
not equal the resource identifier, so it is not used as the reported resource.
With \texttt{enableDetailedLog = true}, the per-resource diagnostic also
reports mission duration, individual cost terms, and excess for candidates
that may subsequently be discarded.

The feasibility loop keeps resource identifiers separate from token counts:
\begin{small}
	\begin{verbatim}
		resource_index=unique(resource_index,'stable');
		resource_index=resource_index(marking_check(resource_index)>=1);
	\end{verbatim}
\end{small}
This handles compatible color matrices containing only some resources.
Transitions already marked as done are skipped without stopping examination
of later transitions. The final token check is retained before a pair is
accepted by the conflict resolver.

\section{Global Completion-Time Criterion}

The local objective $J_c$ and the global completion criterion have different
purposes. The local objective ranks one mission assignment at a time. The
outer simulation evaluates the complete mission list after all batches have
been processed.

A configuration is considered acceptable if its final makespan does not
exceed $8.5$ hours. The interval from $7.5$ to $8.5$ hours is reported as the
target window. A result below $7.5$ hours is also accepted and is classified
as faster than the target. This criterion is implemented separately from the
local cost, so the local controller is not forced to make every intermediate
workload approach eight hours.

The resource assignment history is stored in the batch structure. For each
resource, the simulation records the accumulated time and the mission codes
assigned to it. This allows the final schedule to be evaluated not only by
its makespan, but also by the distribution of missions among the resources.

At the end of each tested resource-count configuration, \texttt{main.m}
computes the actual spread directly from the accumulated workload:
\begin{small}
	\begin{verbatim}
		workloadSpreadHours=(max(batch.Time)-min(batch.Time))/3600;
		balanceStatus="WITHIN TOLERANCE";
		if workloadSpreadHours>balanceOptions.balanceToleranceHours
		balanceStatus="ABOVE TOLERANCE";
		end
	\end{verbatim}
\end{small}
The log prints this value, the configured tolerance, and the status. This
status is diagnostic: the outer search for the number of resources still
accepts configurations according to the existing makespan criterion. It
does not reject a configuration solely because its spread exceeds tolerance.

\section{Prediction Monitoring and Experimental Validation}

For every selected assignment, the predicted workload vector can be stored
in the prediction history. In joint selection mode, only the winning pair
is recorded; predictions of discarded mission candidates are excluded.
The stored prediction can be inspected and compared with the workload after
the CPN update. The current \texttt{plotMpcPrediction.m} displays the terminal
block of the most recent stored prediction as a bar chart in hours; it does
not display a full trajectory. Predicted spread and actual final spread
should be distinguished, because the predictor is invoked during CPN
transition processing and a stored prediction is not itself proof that the
corresponding workload has already been accumulated. Actual final balance
is assessed using \texttt{batch.Time} and its explicit log in \texttt{main.m}.

A suitable validation experiment should compare the following configurations
with identical weights unless weight sensitivity is the quantity being tested:
\begin{enumerate}
	\item resource selection in CPN order, with both flags disabled;
	\item joint mission--resource selection with balance priority disabled;
	\item joint selection with balance priority enabled and a tolerance of
	$0.02$ hours;
	\item additional tolerances or weight combinations for sensitivity analysis.
\end{enumerate}

All strategies must use the same mission list, map, initial resource
positions, and feasibility rules. The comparison should report the mean and
maximum resource workload, the final makespan, the number of deferred
missions, and the computational effort. If several datasets are available,
the weights should be selected using one dataset and tested on another one;
this reduces the risk of choosing weights that are effective only for one
particular Excel file.

The deterministic regression function
\texttt{tests/testGlobalPairSelection.m} uses a temporary path-planning stub
to isolate selection logic from map geometry. It checks joint selection,
skipping completed transitions, restricted resource colors, unavailable
tokens, preservation of CPN order in the legacy mode, and recording only the
winning prediction. The balance checks cover a transition from workloads
$(0.88,0.76)$ hours to $(0.88,0.87)$ hours when the corresponding candidate is
available, rejection of a lower-cost candidate with worse excess, fallback
when every candidate exceeds tolerance, and cost-based selection inside the
tolerance region. These checks are defined in the source; their presence
alone is not evidence of a successful test execution or a full simulation
validation. They can be invoked from the simulation directory with:
\begin{small}
	\begin{verbatim}
		addpath('tests');
		testGlobalPairSelection;
	\end{verbatim}
\end{small}
The illustrative workload values should not be reported as aggregate
experimental results. Quantitative claims about improvement require the
actual dataset, resource count, final workload vector, and execution time
from each comparison run.

\section{Interpretation and Limitations}

The proposed controller is a reduced discrete scheduling MPC. It contains
the essential receding-horizon elements: a state-dependent prediction model,
a finite prediction horizon, a cost criterion, feasibility constraints, and
repeated re-optimization after each state update.

It is not a complete max-plus MPC formulation because the workload dynamics
are represented in conventional algebra and the timed CPN is not converted
into a max-plus state-space model. It is also not a complete MIQP solver,
because the implementation enumerates the currently feasible assignments and
selects the best one instead of optimizing an entire binary sequence over the
horizon.

The main limitations are the following:
\begin{itemize}
	\item the workload model does not explicitly predict vehicle congestion;
	\item the duration depends on A* and therefore on the quality and
	consistency of the map and Excel coordinates;
	\item missions with no feasible path must be deferred or corrected in the
	input data;
	\item the three weights are initially heuristic and require sensitivity
	analysis for a stronger experimental justification;
	\item the future assignment sequence is not optimized jointly over all
	prediction steps;
	\item balance priority can favor a longer assignment or a larger
	overload to reduce spread, and cannot guarantee the final tolerance;
	\item the feasible comparison set remains limited by the current CPN
	place, resource locks, and batch boundaries.
\end{itemize}

These limitations also define possible future developments. A more complete
model could include event times and synchronization through max-plus algebra,
or could formulate the CPN logic and binary assignments as a mixed-logical
dynamical model. Such extensions would allow a full MIQP or discrete-event
MPC optimization, at the cost of increased modeling and computational
complexity.
